\section{Conclusão}

A governança corporativa de Informação e Tecnologia sob a perspectiva do COBIT 2019 consolida-se, ao longo de três décadas de evolução contínua, como um modelo essencial para que organizações contemporâneas transformem investimentos tecnológicos em valor sustentável de negócio. A separação conceitual e prática entre as funções de governança (avaliação, direcionamento e monitoramento pelo conselho e alta liderança) e gestão (planejamento, construção, operação e controle diário por gestores técnicos) estabelece fronteiras nítidas de autoridade, suprimindo sobreposições e eliminando zonas cinzentas de responsabilidade decisória.

Ao incorporar 11 Fatores de Design dinâmicos e métricas de capacidade fundamentadas na escala CMMI, o framework afasta-se definitivamente de receitas padronizadas inflexíveis. Essa arquitetura permite que a governança seja calibrada de acordo com o contexto estratégico, a intensidade regulatória e o apetite ao risco de cada instituição. Ademais, o reconhecimento de que a governança depende da articulação harmônica de sete componentes habilitadores -- e não unicamente de processos -- reforça a importância das estruturas organizacionais colegiadas, dos fluxos informacionais confiáveis e dos aspectos comportamentais e éticos da liderança.

O exame da interoperabilidade do COBIT 2019 comprovou sua vocação como um ``guarda-chuva estratégico'' (meta-modelo). O framework não busca substituir padrões consagrados de mercado, mas atua como o elo integrador superior que orienta o direcionamento e a prestação de contas, delegando a execução tática e operacional para normas consagradas como ITIL (gestão de serviços), ISO/IEC 27001 (segurança da informação), TOGAF (arquitetura corporativa) e PMBOK (projetos).

A análise empírica do estudo de caso no Banco Central da Nigéria (CBN) ilustrou de forma contundente essa dinâmica na prática corporativa: a mera posse de certificações ISO 27001 e processos ITIL não foi suficiente para impedir o caos operacional, a proliferação de \textit{shadow IT} e a aprovação indiscriminada de projetos sem viabilidade técnica. A virada estratégica do CBN ocorreu precisamente com a implantação do COBIT centrada no domínio EDM e na governança compartilhada com diretores executivos de negócio, instituindo triagem unificada de demandas, rastreamento de benefícios do início ao fim e integração dos riscos cibernéticos à Matriz de Risco Corporativo. Em suma, o COBIT 2019 demonstra que a tecnologia só atinge maturidade quando alinhada como ativo nuclear da governança corporativa global.
